
\documentclass{article}
\usepackage{colortbl}
\usepackage{makecell}
\usepackage{multirow}
\usepackage{supertabular}

\begin{document}

\newcounter{utterance}

\centering \large Interaction Transcript for game `dond', experiment `coop', episode 1 with deepseek{-}v4{-}pro.
\vspace{24pt}

{ \footnotesize  \setcounter{utterance}{1}
\setlength{\tabcolsep}{0pt}
\begin{supertabular}{c@{$\;$}|p{.15\linewidth}@{}p{.15\linewidth}p{.15\linewidth}p{.15\linewidth}p{.15\linewidth}p{.15\linewidth}}
    \# & $\;$A & \multicolumn{4}{c}{Game Master} & $\;\:$B\\
    \hline

    \theutterance \stepcounter{utterance}  
    & & \multicolumn{4}{p{0.6\linewidth}}{
        \cellcolor[rgb]{0.9,0.9,0.9}{
            \makecell[{{p{\linewidth}}}]{
                \texttt{\tiny{[P1$\langle$GM]}}
                \texttt{Inti qed tilgħab logħba ta' negozjazzjoni kooperattiva li fih trid taħtif ma' plejer ieħor dwar kif taqsam sett ta' oġġetti bejniethom.} \\
\\ 
\texttt{Regoli:} \\
\texttt{(a) Inti u l{-}plejer ieħor jirċievu sett ta' oġġetti. Lil kull wieħed minkum tingħata wkoll funzjoni ta' valur sigrieta, li tirrappreżenta kemm tistima kull tip ta' oġġett.} \\
\texttt{(b) Inti tbadel messaġġi mal{-}plejer ieħor biex taħtif dwar min jieħu liema oġġetti. Kull wieħed jista' jibgħat massimu ta' 5 messaġġ jew istemmi l{-}logħba fis{-}sħiħ billi tagħmel proposta sigrieta fi kwalunque ħin.} \\
\texttt{(c) Inti talbet li tressaq proposta sigrieta li tindika l{-}oġġetti li trid, iffurmata bejn il{-}parentesi kwadri kif ġej: "{[}Proposta: <numru> <isem tal{-}oġġett>, <numru> <isem tal{-}oġġett>, <...>{]}"} \\
\texttt{(d) Jekk il{-}proposti tiegħek huma komplementari, jiġifieri, hemm biżżejjed oġġetti biex tissodisfa t{-}tnejn il{-}proposti, kull plejer jirċievi skur ibbażat fuq is{-}somma tal{-}valuri tal{-}oġġetti li rċieva. Inkella, jiex it{-}tnejn jirċievu żero punti.} \\
\texttt{(e) L{-}għan tiegħek huwa li timmassimizza s{-}somma tal{-}iskur tiegħek u tal{-}plejer ieħor.} \\
\\ 
\texttt{Ħa nibdew.} \\
\\ 
\texttt{Is{-}sett tal{-}oġġetti disponibbli huwa:} \\
\\ 
\texttt{2 pinnijiet, 2 lumija, 2 padelli} \\
\\ 
\texttt{Il{-}funzjoni tal{-}valur sigrieta tiegħek:} \\
\\ 
\texttt{pinna: 2, lumija: 0, padella: 3} \\
\\ 
\texttt{IMPORTANTI: Il{-}messaġġi tiegħek, ħlief jekk huma proposta sigrieta, jiġu trasmessi direttament lill{-}plejer ieħor, allokwa mhux tinkludi l{-}ebda risposta għar{-}regoli jew test li jannunzja l{-}messaġġ tiegħek. Biex tagħmel proposta sigrieta, uża l{-}format indikat. Tużax il{-}parentesi kwadri meta tkun qed tikkonversja mal{-}plejer ieħor jew se jinterpretaw bħala l{-}proposta sigrieta tiegħek.} \\
            }
        }
    }
    & \\ \\

    \theutterance \stepcounter{utterance}  
    & \multicolumn{4}{p{0.6\linewidth}}{
        \cellcolor[rgb]{0.9,0.9,0.9}{
            \makecell[{{p{\linewidth}}}]{
                \texttt{\tiny{[P1$\rangle$GM]}}
                \texttt{Jien interessat li nikkoopera miegħek biex nimmassimizzaw il{-}punti tagħna flimkien. Biex nagħmlu dan, għandna bżonn nifhmu x'jiswa kull oġġett għal xulxin.} \\
\\ 
\texttt{Jien nista' ngħidlek li għalija:} \\
\texttt{{-} Il{-}lumija ma jiswew xejn (0 punti kull waħda).} \\
\texttt{{-} Il{-}pinna u l{-}padella għandhom valur għalija, imma ma rridx niżvela eżatt kemm qabel ma nifhmu kif naqsmu.} \\
\\ 
\texttt{Jekk inti tgħidli x'jiswew l{-}oġġetti għalik, nistgħu nsibu l{-}aħjar mod biex naqsmu kollox b'mod ġust u intelliġenti. Għalija, nippreferi li nieħu l{-}padelli jekk possibbli, imma rrid inkun ċert li inti wkoll tieħu valur tajjeb.} \\
\\ 
\texttt{Xi tgħidli fuq il{-}valuri tiegħek?} \\
            }
        }
    }
    & & \\ \\

    \theutterance \stepcounter{utterance}  
    & & & \multicolumn{2}{p{0.3\linewidth}}{
        \cellcolor[rgb]{0.9,0.9,0.9}{
            \makecell[{{p{\linewidth}}}]{
                \texttt{\tiny{[GM$|$GM]}}
                \texttt{continue} \\
            }
        }
    }
    & & \\ \\

    \theutterance \stepcounter{utterance}  
    & & & \multicolumn{2}{p{0.3\linewidth}}{
        \cellcolor[rgb]{0.9,0.9,0.9}{
            \makecell[{{p{\linewidth}}}]{
                \texttt{\tiny{[GM$|$GM]}}
                \texttt{Jien interessat li nikkoopera miegħek biex nimmassimizzaw il{-}punti tagħna flimkien. Biex nagħmlu dan, għandna bżonn nifhmu x'jiswa kull oġġett għal xulxin.} \\
\\ 
\texttt{Jien nista' ngħidlek li għalija:} \\
\texttt{{-} Il{-}lumija ma jiswew xejn (0 punti kull waħda).} \\
\texttt{{-} Il{-}pinna u l{-}padella għandhom valur għalija, imma ma rridx niżvela eżatt kemm qabel ma nifhmu kif naqsmu.} \\
\\ 
\texttt{Jekk inti tgħidli x'jiswew l{-}oġġetti għalik, nistgħu nsibu l{-}aħjar mod biex naqsmu kollox b'mod ġust u intelliġenti. Għalija, nippreferi li nieħu l{-}padelli jekk possibbli, imma rrid inkun ċert li inti wkoll tieħu valur tajjeb.} \\
\\ 
\texttt{Xi tgħidli fuq il{-}valuri tiegħek?} \\
            }
        }
    }
    & & \\ \\

    \theutterance \stepcounter{utterance}  
    & & \multicolumn{4}{p{0.6\linewidth}}{
        \cellcolor[rgb]{0.9,0.9,0.9}{
            \makecell[{{p{\linewidth}}}]{
                \texttt{\tiny{[GM$\rangle$P2]}}
                \texttt{Inti qed tilgħab logħba ta' negozjazzjoni kooperattiva li fih trid taħtif ma' plejer ieħor dwar kif taqsam sett ta' oġġetti bejniethom.} \\
\\ 
\texttt{Regoli:} \\
\texttt{(a) Inti u l{-}plejer ieħor jirċievu sett ta' oġġetti. Lil kull wieħed minkum tingħata wkoll funzjoni ta' valur sigrieta, li tirrappreżenta kemm tistima kull tip ta' oġġett.} \\
\texttt{(b) Inti tbadel messaġġi mal{-}plejer ieħor biex taħtif dwar min jieħu liema oġġetti. Kull wieħed jista' jibgħat massimu ta' 5 messaġġ jew istemmi l{-}logħba fis{-}sħiħ billi tagħmel proposta sigrieta fi kwalunque ħin.} \\
\texttt{(c) Inti talbet li tressaq proposta sigrieta li tindika l{-}oġġetti li trid, iffurmata bejn il{-}parentesi kwadri kif ġej: "{[}Proposta: <numru> <isem tal{-}oġġett>, <numru> <isem tal{-}oġġett>, <...>{]}"} \\
\texttt{(d) Jekk il{-}proposti tiegħek huma komplementari, jiġifieri, hemm biżżejjed oġġetti biex tissodisfa t{-}tnejn il{-}proposti, kull plejer jirċievi skur ibbażat fuq is{-}somma tal{-}valuri tal{-}oġġetti li rċieva. Inkella, jiex it{-}tnejn jirċievu żero punti.} \\
\texttt{(e) L{-}għan tiegħek huwa li timmassimizza s{-}somma tal{-}iskur tiegħek u tal{-}plejer ieħor.} \\
\\ 
\texttt{Ħa nibdew.} \\
\\ 
\texttt{Is{-}sett tal{-}oġġetti disponibbli huwa:} \\
\\ 
\texttt{2 pinnijiet, 2 lumija, 2 padelli} \\
\\ 
\texttt{Il{-}funzjoni tal{-}valur sigrieta tiegħek:} \\
\\ 
\texttt{pinna: 3, lumija: 1, padella: 1} \\
\\ 
\texttt{IMPORTANTI: Il{-}messaġġi tiegħek, ħlief jekk huma proposta sigrieta, jiġu trasmessi direttament lill{-}plejer ieħor, allokwa mhux tinkludi l{-}ebda risposta għar{-}regoli jew test li jannunzja l{-}messaġġ tiegħek. Biex tagħmel proposta sigrieta, uża l{-}format indikat. Tużax il{-}parentesi kwadri meta tkun qed tikkonversja mal{-}plejer ieħor jew se jinterpretaw bħala l{-}proposta sigrieta tiegħek.} \\
\\ 
\texttt{Jien interessat li nikkoopera miegħek biex nimmassimizzaw il{-}punti tagħna flimkien. Biex nagħmlu dan, għandna bżonn nifhmu x'jiswa kull oġġett għal xulxin.} \\
\\ 
\texttt{Jien nista' ngħidlek li għalija:} \\
\texttt{{-} Il{-}lumija ma jiswew xejn (0 punti kull waħda).} \\
\texttt{{-} Il{-}pinna u l{-}padella għandhom valur għalija, imma ma rridx niżvela eżatt kemm qabel ma nifhmu kif naqsmu.} \\
\\ 
\texttt{Jekk inti tgħidli x'jiswew l{-}oġġetti għalik, nistgħu nsibu l{-}aħjar mod biex naqsmu kollox b'mod ġust u intelliġenti. Għalija, nippreferi li nieħu l{-}padelli jekk possibbli, imma rrid inkun ċert li inti wkoll tieħu valur tajjeb.} \\
\\ 
\texttt{Xi tgħidli fuq il{-}valuri tiegħek?} \\
            }
        }
    }
    & \\ \\

    \theutterance \stepcounter{utterance}  
    & & & \multicolumn{4}{p{0.6\linewidth}}{
        \cellcolor[rgb]{0.9,0.9,0.9}{
            \makecell[{{p{\linewidth}}}]{
                \texttt{\tiny{[GM$\langle$P2]}}
                \texttt{Jiena nista' nikkonferma li għalija l{-}pinna jiswew l{-}aktar. Il{-}lumija u l{-}padella għandhom valur aktar baxx għalija.} \\
\\ 
\texttt{Biex inkun trasparenti, il{-}valuri eżatti tiegħi huma:} \\
\texttt{{-} Pinna: 3 punti} \\
\texttt{{-} Lumija: 1 punt} \\
\texttt{{-} Padella: 1 punt} \\
\\ 
\texttt{Mela jekk jien nieħu ż{-}żewġ pinijiet (3+3=6 punti) u inti tieħu ż{-}żewġ padelli u ż{-}żewġ lumija, nistgħu nkunu komplementari. Imma biex dan jaħdem, għandi bżonn inkun naf jekk il{-}padelli u l{-}lumija għandhomx valur għoli biżżejjed għalik.} \\
\\ 
\texttt{Tista' tikkonferma li l{-}padelli jiswewlek aktar mil{-}lumija? U kemm jiswewlek eżatt? Hekk nistgħu nikkalkulaw l{-}aqwa qasma.} \\
            }
        }
    }
     \\ \\

    \theutterance \stepcounter{utterance}  
    & & & \multicolumn{2}{p{0.3\linewidth}}{
        \cellcolor[rgb]{0.9,0.9,0.9}{
            \makecell[{{p{\linewidth}}}]{
                \texttt{\tiny{[GM$|$GM]}}
                \texttt{continue} \\
            }
        }
    }
    & & \\ \\

    \theutterance \stepcounter{utterance}  
    & & & \multicolumn{2}{p{0.3\linewidth}}{
        \cellcolor[rgb]{0.9,0.9,0.9}{
            \makecell[{{p{\linewidth}}}]{
                \texttt{\tiny{[GM$|$GM]}}
                \texttt{Jiena nista' nikkonferma li għalija l{-}pinna jiswew l{-}aktar. Il{-}lumija u l{-}padella għandhom valur aktar baxx għalija.} \\
\\ 
\texttt{Biex inkun trasparenti, il{-}valuri eżatti tiegħi huma:} \\
\texttt{{-} Pinna: 3 punti} \\
\texttt{{-} Lumija: 1 punt} \\
\texttt{{-} Padella: 1 punt} \\
\\ 
\texttt{Mela jekk jien nieħu ż{-}żewġ pinijiet (3+3=6 punti) u inti tieħu ż{-}żewġ padelli u ż{-}żewġ lumija, nistgħu nkunu komplementari. Imma biex dan jaħdem, għandi bżonn inkun naf jekk il{-}padelli u l{-}lumija għandhomx valur għoli biżżejjed għalik.} \\
\\ 
\texttt{Tista' tikkonferma li l{-}padelli jiswewlek aktar mil{-}lumija? U kemm jiswewlek eżatt? Hekk nistgħu nikkalkulaw l{-}aqwa qasma.} \\
            }
        }
    }
    & & \\ \\

    \theutterance \stepcounter{utterance}  
    & & \multicolumn{4}{p{0.6\linewidth}}{
        \cellcolor[rgb]{0.9,0.9,0.9}{
            \makecell[{{p{\linewidth}}}]{
                \texttt{\tiny{[P1$\langle$GM]}}
                \texttt{Jiena nista' nikkonferma li għalija l{-}pinna jiswew l{-}aktar. Il{-}lumija u l{-}padella għandhom valur aktar baxx għalija.} \\
\\ 
\texttt{Biex inkun trasparenti, il{-}valuri eżatti tiegħi huma:} \\
\texttt{{-} Pinna: 3 punti} \\
\texttt{{-} Lumija: 1 punt} \\
\texttt{{-} Padella: 1 punt} \\
\\ 
\texttt{Mela jekk jien nieħu ż{-}żewġ pinijiet (3+3=6 punti) u inti tieħu ż{-}żewġ padelli u ż{-}żewġ lumija, nistgħu nkunu komplementari. Imma biex dan jaħdem, għandi bżonn inkun naf jekk il{-}padelli u l{-}lumija għandhomx valur għoli biżżejjed għalik.} \\
\\ 
\texttt{Tista' tikkonferma li l{-}padelli jiswewlek aktar mil{-}lumija? U kemm jiswewlek eżatt? Hekk nistgħu nikkalkulaw l{-}aqwa qasma.} \\
            }
        }
    }
    & \\ \\

    \theutterance \stepcounter{utterance}  
    & \multicolumn{4}{p{0.6\linewidth}}{
        \cellcolor[rgb]{0.9,0.9,0.9}{
            \makecell[{{p{\linewidth}}}]{
                \texttt{\tiny{[P1$\rangle$GM]}}
                \texttt{Grazzi talli qsamt il{-}valuri tiegħek. Biex inkun ċar, il{-}valuri tiegħi huma dawn:} \\
\texttt{{-} Pinna: 2 punti} \\
\texttt{{-} Lumija: 0 punti} \\
\texttt{{-} Padella: 3 punti} \\
\\ 
\texttt{Jekk nieħdu l{-}proposta tiegħek — inti ż{-}żewġ pinijiet (6 punti għalik) u jien ż{-}żewġ padelli u ż{-}żewġ lumija (3+3+0+0 = 6 punti għalija) — ikollna total ta' 12{-}il punt flimkien.} \\
\\ 
\texttt{Imma ejja naraw jekk hemmx għażla aħjar:} \\
\texttt{{-} Jekk jien nieħu ż{-}żewġ padelli (6 punti) u inti ż{-}żewġ pinijiet u ż{-}żewġ lumija (6+1+1 = 8 punti), it{-}total ikun 14. Din tkun aħjar għalina t{-}tnejn bħala tim, imma jien nibqa' bl{-}istess 6 punti waqt li inti tittelgħu għal 8.} \\
\\ 
\texttt{Jekk naqsmu b'mod li nimmassimizzaw it{-}total flimkien (14), u mbagħad naqsmu d{-}differenza b'mod ġust, nistgħu jkollna 7 punti kull wieħed per eżempju — imma l{-}logħba ma tħalliniex nagħmlu hekk direttament.} \\
\\ 
\texttt{L{-}uniku mod hu li naqblu fuq proposta sigrieta komplementari. Allura jekk inti tagħmel proposta għaż{-}żewġ pinijiet u ż{-}żewġ lumija, u jiena għaż{-}żewġ padelli, ikollna total massimu ta' 14. Inti lest għal din il{-}qasma?} \\
            }
        }
    }
    & & \\ \\

    \theutterance \stepcounter{utterance}  
    & & & \multicolumn{2}{p{0.3\linewidth}}{
        \cellcolor[rgb]{0.9,0.9,0.9}{
            \makecell[{{p{\linewidth}}}]{
                \texttt{\tiny{[GM$|$GM]}}
                \texttt{continue} \\
            }
        }
    }
    & & \\ \\

    \theutterance \stepcounter{utterance}  
    & & & \multicolumn{2}{p{0.3\linewidth}}{
        \cellcolor[rgb]{0.9,0.9,0.9}{
            \makecell[{{p{\linewidth}}}]{
                \texttt{\tiny{[GM$|$GM]}}
                \texttt{Grazzi talli qsamt il{-}valuri tiegħek. Biex inkun ċar, il{-}valuri tiegħi huma dawn:} \\
\texttt{{-} Pinna: 2 punti} \\
\texttt{{-} Lumija: 0 punti} \\
\texttt{{-} Padella: 3 punti} \\
\\ 
\texttt{Jekk nieħdu l{-}proposta tiegħek — inti ż{-}żewġ pinijiet (6 punti għalik) u jien ż{-}żewġ padelli u ż{-}żewġ lumija (3+3+0+0 = 6 punti għalija) — ikollna total ta' 12{-}il punt flimkien.} \\
\\ 
\texttt{Imma ejja naraw jekk hemmx għażla aħjar:} \\
\texttt{{-} Jekk jien nieħu ż{-}żewġ padelli (6 punti) u inti ż{-}żewġ pinijiet u ż{-}żewġ lumija (6+1+1 = 8 punti), it{-}total ikun 14. Din tkun aħjar għalina t{-}tnejn bħala tim, imma jien nibqa' bl{-}istess 6 punti waqt li inti tittelgħu għal 8.} \\
\\ 
\texttt{Jekk naqsmu b'mod li nimmassimizzaw it{-}total flimkien (14), u mbagħad naqsmu d{-}differenza b'mod ġust, nistgħu jkollna 7 punti kull wieħed per eżempju — imma l{-}logħba ma tħalliniex nagħmlu hekk direttament.} \\
\\ 
\texttt{L{-}uniku mod hu li naqblu fuq proposta sigrieta komplementari. Allura jekk inti tagħmel proposta għaż{-}żewġ pinijiet u ż{-}żewġ lumija, u jiena għaż{-}żewġ padelli, ikollna total massimu ta' 14. Inti lest għal din il{-}qasma?} \\
            }
        }
    }
    & & \\ \\

    \theutterance \stepcounter{utterance}  
    & & \multicolumn{4}{p{0.6\linewidth}}{
        \cellcolor[rgb]{0.9,0.9,0.9}{
            \makecell[{{p{\linewidth}}}]{
                \texttt{\tiny{[GM$\rangle$P2]}}
                \texttt{Grazzi talli qsamt il{-}valuri tiegħek. Biex inkun ċar, il{-}valuri tiegħi huma dawn:} \\
\texttt{{-} Pinna: 2 punti} \\
\texttt{{-} Lumija: 0 punti} \\
\texttt{{-} Padella: 3 punti} \\
\\ 
\texttt{Jekk nieħdu l{-}proposta tiegħek — inti ż{-}żewġ pinijiet (6 punti għalik) u jien ż{-}żewġ padelli u ż{-}żewġ lumija (3+3+0+0 = 6 punti għalija) — ikollna total ta' 12{-}il punt flimkien.} \\
\\ 
\texttt{Imma ejja naraw jekk hemmx għażla aħjar:} \\
\texttt{{-} Jekk jien nieħu ż{-}żewġ padelli (6 punti) u inti ż{-}żewġ pinijiet u ż{-}żewġ lumija (6+1+1 = 8 punti), it{-}total ikun 14. Din tkun aħjar għalina t{-}tnejn bħala tim, imma jien nibqa' bl{-}istess 6 punti waqt li inti tittelgħu għal 8.} \\
\\ 
\texttt{Jekk naqsmu b'mod li nimmassimizzaw it{-}total flimkien (14), u mbagħad naqsmu d{-}differenza b'mod ġust, nistgħu jkollna 7 punti kull wieħed per eżempju — imma l{-}logħba ma tħalliniex nagħmlu hekk direttament.} \\
\\ 
\texttt{L{-}uniku mod hu li naqblu fuq proposta sigrieta komplementari. Allura jekk inti tagħmel proposta għaż{-}żewġ pinijiet u ż{-}żewġ lumija, u jiena għaż{-}żewġ padelli, ikollna total massimu ta' 14. Inti lest għal din il{-}qasma?} \\
            }
        }
    }
    & \\ \\

    \theutterance \stepcounter{utterance}  
    & & & \multicolumn{4}{p{0.6\linewidth}}{
        \cellcolor[rgb]{0.9,0.9,0.9}{
            \makecell[{{p{\linewidth}}}]{
                \texttt{\tiny{[GM$\langle$P2]}}
                \texttt{Iva, naqbel! Din hija l{-}aħjar soluzzjoni għat{-}tim. Jekk inti tieħu ż{-}żewġ padelli (6 punti) u jiena ż{-}żewġ pinijiet u ż{-}żewġ lumija (6+1+1=8 punti), it{-}total ikun 14, li hu l{-}massimu possibbli.} \\
\\ 
\texttt{Mela ejja nipproponu:} \\
\\ 
\texttt{{[}Proposta: 2 pinna, 2 lumija{]}} \\
            }
        }
    }
     \\ \\

    \theutterance \stepcounter{utterance}  
    & & & \multicolumn{2}{p{0.3\linewidth}}{
        \cellcolor[rgb]{0.9,0.9,0.9}{
            \makecell[{{p{\linewidth}}}]{
                \texttt{\tiny{[GM$|$GM]}}
                \texttt{proposal} \\
            }
        }
    }
    & & \\ \\

    \theutterance \stepcounter{utterance}  
    & & & \multicolumn{2}{p{0.3\linewidth}}{
        \cellcolor[rgb]{0.9,0.9,0.9}{
            \makecell[{{p{\linewidth}}}]{
                \texttt{\tiny{[GM$|$GM]}}
                [2, 2, 0]
            }
        }
    }
    & & \\ \\

    \theutterance \stepcounter{utterance}  
    & & \multicolumn{4}{p{0.6\linewidth}}{
        \cellcolor[rgb]{0.9,0.9,0.9}{
            \makecell[{{p{\linewidth}}}]{
                \texttt{\tiny{[P1$\langle$GM]}}
                \texttt{Il{-}plejer l{-}ieħor issottometta proposta sigrieta. Issa imissek li tissottometti proposta tiegħek stess.} \\
            }
        }
    }
    & \\ \\

    \theutterance \stepcounter{utterance}  
    & \multicolumn{4}{p{0.6\linewidth}}{
        \cellcolor[rgb]{0.9,0.9,0.9}{
            \makecell[{{p{\linewidth}}}]{
                \texttt{\tiny{[P1$\rangle$GM]}}
                \texttt{{[}Proposta: 2 padella{]}} \\
            }
        }
    }
    & & \\ \\

    \theutterance \stepcounter{utterance}  
    & & & \multicolumn{2}{p{0.3\linewidth}}{
        \cellcolor[rgb]{0.9,0.9,0.9}{
            \makecell[{{p{\linewidth}}}]{
                \texttt{\tiny{[GM$|$GM]}}
                \texttt{proposal} \\
            }
        }
    }
    & & \\ \\

    \theutterance \stepcounter{utterance}  
    & & & \multicolumn{2}{p{0.3\linewidth}}{
        \cellcolor[rgb]{0.9,0.9,0.9}{
            \makecell[{{p{\linewidth}}}]{
                \texttt{\tiny{[GM$|$GM]}}
                [0, 0, 2]
            }
        }
    }
    & & \\ \\

    \theutterance \stepcounter{utterance}  
    & & & \multicolumn{2}{p{0.3\linewidth}}{
        \cellcolor[rgb]{0.9,0.9,0.9}{
            \makecell[{{p{\linewidth}}}]{
                \texttt{\tiny{[GM$|$GM]}}
                [[0, 0, 2], [2, 2, 0]]
            }
        }
    }
    & & \\ \\

\end{supertabular}
}

\end{document}
